\documentclass[11pt]{amsart}
\usepackage{amsmath,amssymb,amsthm}
\usepackage[margin=1.1in]{geometry}
\usepackage{xcolor}
\usepackage{hyperref}
\hypersetup{colorlinks=true, linkcolor=blue!60!black, citecolor=blue!60!black, urlcolor=blue!60!black}
% ---------- theorem environments ----------
\theoremstyle{plain}
\newtheorem{theorem}{Theorem}[section]
\newtheorem{proposition}[theorem]{Proposition}
\newtheorem{lemma}[theorem]{Lemma}
\newtheorem{corollary}[theorem]{Corollary}
\newtheorem{conjecture}[theorem]{Conjecture}
\theoremstyle{definition}
\newtheorem{definition}[theorem]{Definition}
\newtheorem{question}[theorem]{Open Question}
\theoremstyle{remark}
\newtheorem{remark}[theorem]{Remark}
% ---------- macros ----------
\newcommand{\R}{\mathbb{R}}
\newcommand{\fpl}[1]{(-\Delta_p)^{s}#1}            % fractional p-Laplacian
\newcommand{\Jp}[1]{|#1|^{p-2}#1}                  % J_p(t) = |t|^{p-2} t
\newcommand{\al}{\alpha}
\newcommand{\Tail}{\mathrm{Tail}_{p-1,sp}}
\newcommand{\TODO}[1]{\textcolor{red}{\textbf{TODO:} #1}}
\newcommand{\NOTE}[1]{\textcolor{blue!70!black}{\textit{Note: #1}}}
\title[Notes: nonlocal $p$-dead-core]{Working notes: improved regularity for a nonlocal\\ two-phase dead-core problem driven by the fractional $p$-Laplacian}
\author{Hikmatullo}
\date{\today\ --- lab notebook, not a draft}
\begin{document}
\maketitle
\tableofcontents
%======================================================================
\section{Introduction (draft)}\label{sec:intro}
%======================================================================
Dead-core phenomena arise when a diffusing reactant is consumed by a
sufficiently strong absorption: the reactant vanishes identically on an
interior region, the \emph{dead core}, on whose boundary the solution
touches zero together with its gradient. The defining mathematical feature
is a sublinear, non-Lipschitz absorption $u^\gamma$ with $0<\gamma<1$
(local) or $0<\gamma<p-1$ (the $p$-Laplacian scaling): it is precisely the
failure of Lipschitz continuity at $u=0$ that permits a solution to reach
zero on a set of positive measure rather than at isolated points. We study
the \emph{two-phase, nonlocal} model
\begin{equation}\label{eq:intro}
  \fpl{u}=u_+^{\gamma}-u_-^{\gamma}\qquad\text{in }B_1,
  \qquad 0<\gamma<p-1,\ s\in(0,1),\ p\ge2,
\end{equation}
where $\fpl{}$ is the fractional $p$-Laplacian. Two structural choices
distinguish \eqref{eq:intro} from the classical literature. First, the
problem is \emph{nonlocal}: the value of $\fpl{u}(x)$ depends on $u$ over
all of $\R^n$, so the maximum-principle and minimum-of-the-energy arguments
that drive the local one-phase theory are unavailable --- in the nonlocal
setting a comparison principle requires sign information on the entire
complement of $B_1$. Second, the problem is \emph{two-phase}: $u$ is not
sign-constrained, so branching points are not local extrema and the
one-phase tools fail for a second, independent reason. The one-phase
signed problem is in fact obstructed nonlocally by the strong minimum
principle, which makes the two-phase formulation the natural setting.

Our main result is a sharp growth estimate at branching points. At a point
$x_0$ with $u(x_0)=|Du(x_0)|=0$,
\begin{equation}\label{eq:introgrowth}
  |u(x)|\le C\,\|u\|_{L^\infty(\R^n)}\,|x-x_0|^{\al},
  \qquad \al=\frac{sp}{p-1-\gamma},
\end{equation}
an exponent strictly larger than the optimal regularity
$\min\!\big(1,\tfrac{sp}{p-1}\big)$ of $(s,p)$-harmonic functions, since
$\gamma>0$: solutions are \emph{more} regular at branching points than the
homogeneous theory predicts. Passing to the limit $s\nearrow1$ yields the
analogous estimate for the local two-phase $p$-dead-core problem
$\Delta_pu=u_+^\gamma-u_-^\gamma$ with exponent $\frac{p}{p-1-\gamma}$,
which appears to be new even locally --- the existing $p$-dead-core results
(da Silva--Salort, da Silva--Rossi--Salort, Teixeira) are one-phase.

The method adapts the blow-up/Liouville strategy of dos
Prazeres--Teymurazyan--Urbano \cite{PTU} (the case $p=2$) and its
quasilinear local counterpart in Correa--dos Prazeres \cite{CP}. One
rescales according to the natural exponent $\al$, shows the absorption term
vanishes along the blow-up, and classifies the limit by a Liouville-type
theorem. The new mechanism for $p\neq2$ is that the \emph{$(p-1)$-homogeneity}
of $\fpl{}$ replaces the linearity used in \cite{PTU}: it is what forces the
right-hand side to vanish in the limit (Proposition~\ref{prop:homog}),
playing the role that differentiability of the nonlinearity at $0$ plays in
the fully nonlinear case. The principal difficulties are genuinely
nonlocal: controlling the far-field tail of the blow-up family uniformly
(\S\ref{sec:constraints}), and proving the Liouville theorem without the
linear structure (\S\ref{sec:liouville}); the interior $C^{1,\al_0}$ input
is the recent estimate of Giovagnoli--Jesus--Silvestre \cite{GJS}.

\NOTE{Intro drafted 2026-06-15. Keep the exponent claim in
\eqref{eq:introgrowth} phrased against ``optimal regularity of
$(s,p)$-harmonic functions''; the precise relation of $\al$ to $1+\al_0$
is the open bookkeeping point flagged in \S\ref{sec:targets}.}
%======================================================================
\section{The problem and the program}
%======================================================================
We study the two-phase nonlocal dead-core problem
\begin{equation}\label{eq:main}
  \fpl{u} = u_+^{\gamma} - u_-^{\gamma} \qquad \text{in } B_1,
\end{equation}
with $0<\gamma<p-1$, $s\in(0,1)$, $p\ge 2$, where
\[
  \fpl{u}(x) := c_{n,s,p}\;\mathrm{P.V.}\!\int_{\R^n}
    \frac{\Jp{\big(u(x)-u(y)\big)}}{|x-y|^{n+sp}}\,dy .
\]
The goal is the analogue of the Prazeres--Teymurazyan--Urbano (PTU) result
\cite{PTU}: at free boundary points $x_0$ with $u(x_0)=|Du(x_0)|=0$,
\[
  |u(x)| \le C\,\|u\|_{L^\infty(\R^n)}\,|x-x_0|^{\al},
  \qquad
  \al := \frac{sp}{\,p-1-\gamma\,},
\]
which exceeds the optimal regularity $\min\!\big(1,\tfrac{sp}{p-1}\big)$ of
$(s,p)$-harmonic functions, since $\gamma>0$.
\medskip
\noindent\textbf{Program (the PTU machine, transplanted).}
\begin{enumerate}
  \item[\S\ref{sec:scaling}] Scaling exponent (Gear 2). \emph{Survives.}
  \item[\S\ref{sec:homog}] Blow-up limit equation via $(p-1)$-homogeneity
        (replaces linearization). \emph{Survives, new mechanism.}
  \item[\S\ref{sec:constraints}] Constraint region in $(s,p,\gamma)$
        (tail integrability, $\beta<1$, regularity ceiling).
  \item[\S\ref{sec:liouville}] Liouville theorem for entire
        $(s,p)$-harmonic functions with controlled $C^{1,\al_0}$ growth.
        \emph{The load-bearing wall; likely the main new lemma.}
  \item[\S\ref{sec:targets}] Target theorems: growth at the free boundary,
        gradient estimate, Liouville with growth, the limit $s\nearrow 1$
        (two-phase local $p$-dead-core: new even locally).
\end{enumerate}
\NOTE{One-phase signed dead cores are obstructed nonlocally by the strong
minimum principle; the two-phase formulation \eqref{eq:main} is the natural
(and essentially only) nontrivial setting. Use this in the introduction.}
%======================================================================
\section{Notation and standing assumptions}
%======================================================================
For $\al\in(0,1]$ we use the seminorms $[\,\cdot\,]_{C^{\al}(B_R)}$ and
$[\,\cdot\,]_{C^{1+\al}(B_R)}$ as in \cite[Sec.~2]{PTU}.
The nonlocal tail adapted to $\fpl{}$ is
\[
  \Tail(u;x_0,R)
  := \left( R^{sp} \int_{\R^n\setminus B_R(x_0)}
       \frac{|u(y)|^{p-1}}{|y-x_0|^{n+sp}}\,dy \right)^{\!\frac{1}{p-1}},
\]
replacing the space $L^1_s(\R^n)$ of the case $p=2$.
\TODO{Fix the notion of solution for the first pass: \emph{weak} solutions
$u\in W^{s,p}_{loc}\cap L^\infty(\R^n)$, with the viscosity equivalence
(Korvenp\"a\"a--Kuusi--Lindgren) cited later. Existence plumbing deferred
to the final write-up.}
%======================================================================
\section{Gear 2: the scaling exponent}\label{sec:scaling}
%======================================================================
\begin{proposition}[Scaling]\label{prop:scaling}
Let $u$ solve \eqref{eq:main} and, for $r>0$ and $\al>0$, set
$u_r(x):=u(rx)/r^{\al}$. Then
\[
  \fpl{u_r}(x) \;=\; r^{\,sp-\al(p-1)}\,\big(\fpl{u}\big)(rx),
\]
and $u_r$ solves \eqref{eq:main} (in $B_{1/r}$) if and only if
\[
  sp-\al(p-1)+\al\gamma=0,
  \qquad\text{i.e.}\qquad
  \boxed{\;\al=\frac{sp}{p-1-\gamma}\;}.
\]
\end{proposition}
\begin{proof}
Write $J_p(t):=\Jp{t}$ and note the $(p-1)$-homogeneity
$J_p(\lambda t)=\lambda^{p-1}J_p(t)$ for $\lambda>0$.
Fix $x\in B_{1/r}$. Since $u_r(x)-u_r(y)=r^{-\al}\big(u(rx)-u(ry)\big)$,
\[
  \fpl{u_r}(x)
  = c_{n,s,p}\,r^{-\al(p-1)}\,\mathrm{P.V.}\!\int_{\R^n}
    \frac{J_p\big(u(rx)-u(ry)\big)}{|x-y|^{n+sp}}\,dy .
\]
Change variables $z=ry$, so $dy=r^{-n}dz$ and
$|x-y|^{-(n+sp)}=r^{\,n+sp}\,|rx-z|^{-(n+sp)}$; the kernel and the volume
element together contribute $r^{sp}$, and we obtain
\begin{equation}\label{eq:scaling-identity}
  \fpl{u_r}(x)
  = r^{\,sp-\al(p-1)}\,
    c_{n,s,p}\,\mathrm{P.V.}\!\int_{\R^n}
    \frac{J_p\big(u(rx)-u(z)\big)}{|rx-z|^{n+sp}}\,dz
  = r^{\,sp-\al(p-1)}\,\big(\fpl{u}\big)(rx).
\end{equation}
Now use the equation. Since $r^{-\al}>0$, positive and negative parts
scale as $(u_r)_\pm(x)=r^{-\al}u_\pm(rx)$, whence
\[
  (u_r)_+^{\gamma}(x)-(u_r)_-^{\gamma}(x)
  = r^{-\al\gamma}\Big(u_+^{\gamma}-u_-^{\gamma}\Big)(rx).
\]
Combining this with \eqref{eq:scaling-identity} and
$\fpl{u}(rx)=\big(u_+^{\gamma}-u_-^{\gamma}\big)(rx)$ for $rx\in B_1$,
\[
  \fpl{u_r}(x)
  = r^{\,sp-\al(p-1)+\al\gamma}\,
    \Big[(u_r)_+^{\gamma}-(u_r)_-^{\gamma}\Big](x),
  \qquad x\in B_{1/r}.
\]
Hence $u_r$ solves \eqref{eq:main} in $B_{1/r}$ if and only if
$sp-\al(p-1)+\al\gamma=0$, i.e.\ $\al=\frac{sp}{p-1-\gamma}$
(recall $\gamma<p-1$, so the denominator is positive).

The same identity holds for weak solutions: for
$\varphi\in C_c^\infty(B_{1/r})$, the change of variables
$(x,y)\mapsto(rx,ry)$ in the bilinear form
\[
  \mathcal{E}(u,\varphi)
  := \frac{c_{n,s,p}}{2}\iint_{\R^n\times\R^n}
     \frac{J_p\big(u(x)-u(y)\big)\big(\varphi(x)-\varphi(y)\big)}
          {|x-y|^{n+sp}}\,dx\,dy
\]
gives $\mathcal{E}(u_r,\varphi)=r^{\,sp-\al(p-1)-n}\,
\mathcal{E}\big(u,\varphi(\cdot/r)\big)$, and the right-hand side scales
by the matching factor $r^{-\al\gamma-n}$ after the same substitution;
the exponent balance is identical.

\emph{Consistency checks.} For $p=2$:
$\al=\frac{2s}{1-\gamma}$, the PTU exponent \cite{PTU}. For $s\nearrow1$
(with $p$ fixed): $\al\to\frac{p}{p-1-\gamma}$, the local dead-core
exponent of da Silva--Salort. Finally, for $\gamma>0$ one has
$\al=\frac{sp}{p-1-\gamma}>\frac{sp}{p-1}$, the optimal exponent of
$(s,p)$-harmonic functions, as claimed.
\end{proof}
\begin{remark}[Homogeneity in the dependent variable]\label{rem:lambda}
Alongside \eqref{eq:scaling-identity} we record, for $\lambda>0$,
$\fpl{(\lambda u)}=\lambda^{p-1}\,\fpl{u}$, which is the identity used in
the blow-up argument of \S\ref{sec:homog}.
\end{remark}
\begin{remark}
$\al>\frac{sp}{p-1}$ for every $\gamma>0$: the target exponent strictly
beats the optimal regularity of the homogeneous equation. This is the
``better than Schauder'' punchline of the paper.
\end{remark}
%======================================================================
\section{Gear 4$'$: the blow-up limit equation via homogeneity}\label{sec:homog}
%======================================================================
In PTU, linearity makes incremental quotients $s$-harmonic. For $p\neq2$
this fails, but $(p-1)$-homogeneity of $\fpl{}$ makes linearization
unnecessary: the absorption term dies in the blow-up.
\begin{proposition}[Vanishing right-hand side along the blow-up]\label{prop:homog}
Let $u_k$ be normalized solutions of \eqref{eq:main}, let $r_k\to0$,
$\theta_k\to\infty$, and set
\[
  v_k(x):=\frac{u_k(r_kx)}{\theta_k\, r_k^{\al}},
  \qquad \al=\frac{sp}{p-1-\gamma}.
\]
Then
\[
  \fpl{v_k}(x)
  = \theta_k^{\,\gamma-(p-1)}
    \Big[ (v_k)_+^{\gamma}(x) - (v_k)_-^{\gamma}(x) \Big],
\]
and since $\gamma<p-1$, the right-hand side tends to $0$ locally uniformly
whenever $v_k$ is locally bounded. Consequently any local uniform limit
$v_0$ is an entire $(s,p)$-harmonic function:
$\fpl{v_0}=0$ in $\R^n$.
\end{proposition}
\begin{proof}
Write $\lambda_k:=(\theta_k r_k^{\al})^{-1}>0$, so that
$v_k(x)=\lambda_k\,u_k(r_kx)$. We combine two scalings.

\emph{Spatial scaling.} For the dilation $x\mapsto r_kx$, the change of
variables of Proposition~\ref{prop:scaling} (with the multiplicative
homogeneity $\fpl{(\lambda w)}=\Jp{\lambda}\,\fpl{w}=\lambda^{p-1}\fpl{w}$
of Remark~\ref{rem:lambda} absorbing the prefactor) gives the order-$sp$
identity
\[
  \fpl{\big[u_k(r_k\,\cdot)\big]}(x)
  = r_k^{\,sp}\,\big(\fpl{u_k}\big)(r_kx).
\]

\emph{Amplitude scaling.} By $(p-1)$-homogeneity in the dependent variable,
\[
  \fpl{v_k}(x)
  = \fpl{\big[\lambda_k\,u_k(r_k\,\cdot)\big]}(x)
  = \lambda_k^{\,p-1}\,r_k^{\,sp}\,\big(\fpl{u_k}\big)(r_kx).
\]

\emph{Use the equation.} For $r_kx\in B_1$, $\fpl{u_k}(r_kx)
=u_k(r_kx)_+^{\gamma}-u_k(r_kx)_-^{\gamma}$. Since $\lambda_k>0$,
$(v_k)_\pm(x)=\lambda_k\,u_k(r_kx)_\pm$, whence
$u_k(r_kx)_\pm^{\gamma}=\lambda_k^{-\gamma}(v_k)_\pm^{\gamma}(x)$ and
\[
  \fpl{v_k}(x)
  = \lambda_k^{\,p-1-\gamma}\,r_k^{\,sp}
    \Big[(v_k)_+^{\gamma}-(v_k)_-^{\gamma}\Big](x).
\]

\emph{The exponents collapse.} By definition of $\lambda_k$,
$\lambda_k^{\,p-1-\gamma}r_k^{\,sp}
=\theta_k^{-(p-1-\gamma)}\,r_k^{\,sp-\al(p-1-\gamma)}$, and
$\al(p-1-\gamma)=sp$ \emph{exactly} (this is the defining property of
$\al$), so the power of $r_k$ vanishes and
\[
  \fpl{v_k}(x)
  = \theta_k^{\,\gamma-(p-1)}\Big[(v_k)_+^{\gamma}-(v_k)_-^{\gamma}\Big](x),
  \qquad x\in B_{1/r_k}.
\]
Since $0<\gamma<p-1$, the exponent $\gamma-(p-1)<0$ and $\theta_k\to\infty$,
so the prefactor $\to0$; if $v_k$ is locally bounded the bracket is locally
bounded and the right-hand side $\to0$ locally uniformly. Passing to the
limit (stability of weak/viscosity solutions of $\fpl{}$ under local
uniform convergence with tail control, see the discussion below) yields
$\fpl{v_0}=0$ in $\R^n$.
\end{proof}
\NOTE{The cancellation is the whole mechanism: spatial scaling supplies
$r_k^{sp}$, amplitude scaling supplies $(\theta_k r_k^{\al})^{-(p-1-\gamma)}$,
and the choice $\al=\frac{sp}{p-1-\gamma}$ is precisely what makes the two
powers of $r_k$ cancel, leaving the pure damping factor
$\theta_k^{-(p-1-\gamma)}\to0$. This is what we mean by ``homogeneity
replaces linearization''. The only analytic input still to pin down is the
stability lemma for $\fpl{}$ (Brasco--Lindgren / Korvenp\"a\"a--Kuusi--Lindgren),
the $p\neq2$ replacement for \cite[Lemma 5]{CS-approx}; it carries the same
\emph{global} tail-convergence hypothesis whose verification is the subject
of \S\ref{sec:constraints}.}
\NOTE{This parallels PTU \S6 (fully nonlinear case), where
$\delta^{-1}F(\delta M)\to DF(0)M$; here homogeneity plays the role of
differentiability of $F$ at $0$ --- which $\fpl{}$ does not have.
Worth a remark in the paper.}
%======================================================================
\section{Gear 3 bookkeeping: the admissible region in $(s,p,\gamma)$}\label{sec:constraints}
%======================================================================
Three constraints must hold simultaneously. Set $\al=\frac{sp}{p-1-\gamma}$
and $\beta:=\al-1$ (the $\nu=1$ case).
\begin{lemma}[Tail integrability]\label{lem:tail}
The growth $|v_0(x)|\lesssim |x|^{\al}$ at infinity is compatible with
$\Tail(v_0;0,R)<\infty$ for all $R$ if and only if $\al(p-1)<sp+p-1$,
i.e.
\[
  \boxed{\;\gamma \;<\; (p-1)-\frac{sp}{1+\,sp/(p-1)}
  \;=\;(p-1)\Big(1-\frac{sp}{p-1+sp}\Big)
  \;=\;\frac{(p-1)^2}{p-1+sp}\;}
\]
\end{lemma}
\begin{proof}[Verification (2026-06-10; constants to be rechecked)]
We record the bookkeeping for the blow-up family
$v_k(x)=\frac{u_k(r_kx)}{\theta_k r_k^{\al}}$, $\|u_k\|_{L^\infty(\R^n)}\le M$,
under the seminorm normalization
$\theta_k=\theta_k(r_k)$ with
$\theta_k(r'):=\sup_{r'<r<1/2} r^{1+\al'-\al}[u_k]_{C^{1+\al'}(B_r)}$.

\emph{Step 1: uniform tails for $v_k$ itself are impossible.}
The interior information gives the growth bound $|v_k(y)|\le C|y|^{\al}$
for $1\le|y|\le R^*_k:=\tfrac{1}{2r_k}$ (gradient-anchor trick of
\cite{CP}), while globally $|v_k|\le M(\theta_k r_k^{\al})^{-1}$.
The two bounds cross at
$\rho_k=(M/C\theta_k)^{1/\al}\,r_k^{-1}$. Splitting the tail there and
using $\al(p-1)-sp=\al\gamma>0$,
\[
  \Tail^{\,p-1}(v_k;0,R)
  \;\le\; c\,R^{\al(p-1)}
  \;+\; c\,R^{sp}\,\rho_k^{\al\gamma}
  \;=\; c\,R^{\al(p-1)} + c\,R^{sp}\big(\theta_k r_k^{\al}\big)^{-\gamma},
\]
with both the growth region and the capped region contributing
comparably. But the normalization forces
$\theta_k\le [u_k]_{C^{1+\al'}(B_{1/2})}\, r_k^{1+\al'-\al}\le
2C\,r_k^{1+\al'-\al}$, i.e.\ $\theta_k r_k^{\al}\le 2C\,r_k^{1+\al'}\to0$;
hence $\sup_k\Tail(v_k;0,R)=\infty$. \textbf{Conclusion: the
$C^{1,\al_0}$ estimate can never be applied to $v_k$ directly; one must
work one order down.} This settles the ``decide'' in the old TODO: the
incremental-quotient step (PTU (4.11)) is \emph{not} optional.

\emph{Step 2: incremental quotients give exactly the boxed condition.}
For $w_k^h:=\frac{v_k(\cdot+h)-v_k(\cdot)}{|h|}$, $0<|h|\le1$, the
gradient bound yields $|w_k^h(y)|\le C(1+|y|)^{\al-1}$ for
$|y|\le R^*_k$, and now the growth-region tail integral
\[
  R^{sp}\int_{|y|>R}|y|^{(\al-1)(p-1)-n-sp}\,dy
\]
converges \emph{uniformly in $k$} if and only if $(\al-1)(p-1)<sp$,
which upon substituting $\al=\frac{sp}{p-1-\gamma}$ is precisely
\[
  \gamma<\frac{(p-1)^2}{p-1+sp}.
\]
Sanity check at $p=2$: $\gamma<\frac{1}{1+2s}$, matching PTU, and
$\gamma<1/3$ uniformly for $s\in(\frac12,1)$. \checkmark

\emph{Step 3: the residual far-field gap.}
Beyond $R^*_k$ the cap gives the contribution
$c\,R^{sp}|h|^{-(p-1)}\,\theta_k^{-(p-1)}r_k^{-\al\gamma}$, which is
bounded iff $\theta_k\ge c\,r_k^{-\al\gamma/(p-1)}$ --- a \emph{lower}
bound on $\theta_k$ that the contradiction hypothesis does not obviously
provide ($\theta_k\to\infty$, but at an unknown rate relative to $r_k$).
\TODO{(Checked 2026-06-10.) PTU \S4 does \emph{not} obviously provide
this bound: their (4.11) rests on (4.8), valid only for
$|x|\le\frac{1}{2r_k}$, while the limit passage (4.12) needs uniform
$L^1_s$ control on all of $\R^n$; the outer region contributes
$\sim\theta_k^{-1}r_k^{-\frac{2s\gamma}{1-\gamma}}$ --- the same lower
bound, apparently unaddressed. So the far-field issue is not an artifact
of $p\neq2$: it is present (implicitly resolved, or genuinely open) in
the template itself. \textbf{Priority question for R.T.} Candidate
fixes: dichotomy $\theta_k\gtrless r_k^{-\delta}$; cutoff with error
estimate. (Checked: weakening the stability lemma is \emph{not} an
option --- \cite{CS-approx} Lemma 4.3 requires global $L^1(\omega)$
convergence, and their Remark 4.4 exhibits a counterexample of exactly
this far-field-mass type.)}
\end{proof}
\begin{lemma}[Order of the jet: $\beta<1$]\label{lem:beta}
The Liouville theorem with $\nu=1$ requires
$\beta=\frac{sp}{p-1-\gamma}-1<1$, i.e.
\[
  \boxed{\;\gamma<(p-1)-\frac{sp}{2}\;}
\]
which also forces $sp<2(p-1)$.
\end{lemma}
\begin{lemma}[Regularity ceiling]\label{lem:ceiling}
The compactness step requires a uniform interior estimate one derivative
above the vanishing jet. With $\nu=1$ this means local $C^{1,\al_0}$
estimates for normalized solutions of $\fpl{u}=f\in L^\infty$, currently
available for
\[
  p\in\Big[\,2,\ \tfrac{2}{1-s}\,\Big),
\]
together with the constraint $\al<1+\al_0$.
\end{lemma}
\NOTE{(2026-06-10, after reading \cite{GJS}.) Two structural facts.
(i) \cite{GJS} is stated for $\fpl{u}=0$; by their own remark the
extension to $f\in L^\infty$ is expected when the linearized operator
has order $2-p(1-s)>1$, i.e.\ $p<\frac{1}{1-s}$ --- which forces
$s>\frac12$, matching PTU. Their stage-3 machinery (Imbert--Silvestre)
already takes a right-hand side, and the Lipschitz starting point with
data exists (Biswas--Sen, arXiv:2507.09920). Writing this extension out
is a candidate lemma of the present paper.
(ii) \textbf{Warning for Theorem~\ref{thm:local}:} the constants in
\cite{GJS} are explicitly NOT uniform as $s\nearrow1$ (their \S1.2:
the first stage is ``purely nonlocal''). The $s\nearrow1$ route to the
local two-phase theorem is blocked through this estimate; the local
theorem should instead be proved directly (classical $C^{1,\alpha}$
theory for the $p$-Laplacian, run through the scheme of \cite{CP} \S6
with right-hand side $u_+^\gamma-u_-^\gamma$).}
\TODO{(i) Track down the $C^{1,\al_0}$ paper (cited as [28] in
arXiv:2509.15988) and record the exact statement: dependence of $\al_0$
and of the constant on $n,s,p$, presence of tail terms, and uniformity as
$s\nearrow1$ (needed for \S\ref{sec:targets}, Theorem~\ref{thm:local}).
(ii) Intersect the three boxed regions and draw the admissible set; this
intersection \emph{is} the hypothesis of the main theorem, the analogue of
$\gamma\in(0,1/3)$, $s\in(1/2,1)$ in PTU.
(iii) Decide whether a weaker fallback (only $C^{0,\gamma_0}$ estimates,
$\gamma_0=\min(1,\frac{sp}{p-1})$, with the free boundary condition
$u(x_0)=0$ only and $\al\le\gamma_0$ forced) yields a nonempty parameter
range --- this would give an unconditional but weaker theorem.}
\subsection*{Attack notes on the far-field bound (2026-06-10; heuristic,
exponents to be redone slowly on paper)}
Write the needed bound as
\[
  \textbf{(N)}\qquad \theta_k^{\,p-1}\, r_k^{\,\al\gamma}\;\ge\; c\;>\;0 .
\]

\emph{(a) Compatibility constraint on the working seminorm order.}
The device forces $\theta_k\le 2C\,r_k^{1+\al'-\al}$, so (N) is
consistent with the worst case only if
$r^{1+\al'-\al}\gtrsim r^{-\al\gamma/(p-1)}$ for small $r$, i.e.
\[
  1+\al' \;\le\; \al\,\frac{p-1-\gamma}{p-1} \;=\; \frac{sp}{p-1},
  \qquad\text{equivalently}\qquad
  \al' \le \frac{sp-(p-1)}{p-1},
\]
which requires $s>1-\frac1p$. At $p=2$ this reads $s>\frac12$ ---
precisely the PTU range; plausibly the structural reason for it.

\emph{(b) Far field as a slowly varying datum --- the likely implicit
fix in PTU.} For $x\in B_R$, split the operator at $R^*_k=\frac{1}{2r_k}$
and treat the far region as a datum $F_k(x)$. Its \emph{size} diverges
like $\theta_k^{-(p-1)}r_k^{-\al\gamma}$ when (N) fails --- but its
\emph{oscillation} on $B_R$ is much better. At $p=2$: the divergent part
of $F_k$ is constant in $x$ and increments kill it; the $x$-variation is
controlled by one extra power of kernel decay,
\[
  \mathrm{cap}_k\,(R^*_k)^{-2s-1}
  = M\theta_k^{-1}r_k^{\,2s+1-\al}\longrightarrow 0
  \iff \gamma<\tfrac{1}{1+2s},
\]
\emph{exactly the tail condition.} So for $p=2$ the limit equation for
increments survives without any lower bound on $\theta_k$: stability
should be applied to the \emph{localized} function with the far field
carried as a datum of vanishing oscillation, not via global $L^1_s$
convergence. \TODO{Make this rigorous at $p=2$ first; it would repair
PTU (4.12) and is publishable as a remark/appendix.}

\emph{(c) The $p\neq2$ version.} Truncate: $\tilde v_k=v_k\chi_{B_{R^*_k/2}}$
solves the equation in $B_R$ up to an error
$G_k(x)=c\int_{|y|>R^*_k/2}[\Jp{(v_k(x)-v_k(y))}-\Jp{v_k(x)}]\,|x-y|^{-n-sp}dy$.
Decompose $G_k=a_k+g_k$, $a_k$ constant. Two estimates:
\[
  \|\nabla g_k\|_{L^\infty(B_R)}
  \lesssim
  \underbrace{\mathrm{cap}_k^{\,p-2}(R^*_k)^{-sp}}_{\to\,0\ \text{since}\ sp-\al(p-2)=\al(1-\gamma)>0}
  +\;
  \underbrace{\mathrm{cap}_k^{\,p-1}(R^*_k)^{-sp-1}}_{\to\,0\ \iff\ \gamma<\frac{p-1}{1+sp}} .
\]
The new admissibility condition
\[
  \boxed{\;\gamma<\frac{p-1}{1+sp}\;}
\]
coincides with the tail condition $\gamma<\frac{(p-1)^2}{p-1+sp}$ at
$p=2$ and is \emph{stricter} for $p>2$; if (b)/(c) are the right
mechanism, this replaces (or joins) Lemma~\ref{lem:tail} in the
hypotheses. The divergent constant $a_k$: for $p=2$ increments kill it;
for $p\neq2$ \TODO{decide what kills $a_k$ --- a posteriori boundedness
(if $a_k\to\infty$, interior estimates on the now-tail-finite
$\tilde v_k$ force a contradiction in the equation), or a smarter
truncation. This is the crux for $p\neq2$.}

\emph{(c$'$) Proposed repair at $p=2$ (2026-06-10; complete in outline,
every step to be checked).} Let $w_k=v_k(\cdot+h)-v_k(\cdot)$ as in PTU
(4.11)--(4.12), $\rho_k:=\frac{1}{4r_k}$, and truncate:
$\tilde w_k:=w_k\chi_{B_{\rho_k}}$. For $x\in B_R$,
\[
  (-\Delta)^s\tilde w_k(x) = f_k(x) + H_k(x),
  \qquad
  H_k(x):=c\int_{|y|>\rho_k} w_k(y)\,|x-y|^{-n-2s}\,dy,
\]
with $f_k\to0$ locally uniformly (PTU (4.12)). Then:
\begin{enumerate}
\item \emph{Oscillation of $H_k$ vanishes:} $|\nabla H_k|\lesssim
  C_h\rho_k^{\al-2s-2} + \theta_k^{-1}r_k^{\,2s+1-\al}\to0$, using the
  growth bound in the annulus and the cap beyond, with
  $2s+1-\al>0\iff\gamma<\frac{1}{1+2s}$. So $H_k=a_k+o(1)$ with $a_k$
  constant.
\item \emph{$a_k$ is bounded a posteriori:} $\tilde w_k$ has uniform
  local $C^{2s+\gamma'}$ bounds and uniform $L^1_s$ norm (growth
  $|y|^{\al-1}$, integrable by the tail condition), so
  $|(-\Delta)^s\tilde w_k(x_0)|\le A$ uniformly at a fixed point; since
  $f_k\to0$, $|a_k|\le A+o(1)$. (No lower bound on $\theta_k$ ever
  needed.)
\item \emph{Limit equation with a constant:} dominated convergence gives
  $\tilde w_k\to w_0$ in $L^1_s$; CS stability now applies and yields
  $(-\Delta)^s\big(v_0(\cdot+h)-v_0\big)=a(h)$ in $B_R$, all $R$, where
  $a(h):=\lim a_k$ (subsequence).
\item \emph{$a(h)$ is linear in $h$:} the cocycle identity
  $w^{h_1+h_2}=w^{h_1}(\cdot+h_2)+w^{h_2}$ and translation invariance
  give $a(h_1+h_2)=a(h_1)+a(h_2)$; with measurability, $a(h)=b\cdot h$.
\item \emph{Second differences kill $b$:}
  $\Delta_h^2v_0 = w^{h}(\cdot+h)-w^{h}$ satisfies
  $(-\Delta)^s\Delta^2_hv_0 = a(h)-a(h)=0$ entire, with growth
  $|\Delta^2_hv_0(x)|\le C_h(1+|x|)^{\al-2}$. PTU Theorem~3.1
  ($\nu=1$, $\beta=\al-2<1$) applies: $\Delta^2_hv_0$ is affine for
  every $h$, hence $v_0$ is a polynomial of degree $\le3$; growth
  $\al<1+2s<3$ cuts it to degree $\le2$; the vanishing $2$-jet at the
  branching point ($\nu=2$) gives $v_0\equiv0$, contradicting the
  normalization. \emph{The Liouville input is even weaker than PTU's
  original use.}
\end{enumerate}
\TODO{Check (1)--(5) line by line, especially: the uniform local
regularity needed in (2) (PTU's $C^{2s+\gamma}$ bounds suffice?);
measurability/continuity of $a(h)$ in (4); and that the degree-$\le3$
deduction in (5) is correct ($\Delta^2_h v$ affine $\forall h$ $\iff$
$v$ cubic --- standard, but verify the nonsmooth case). If this holds,
it repairs PTU (4.12) \emph{and} removes the need for any far-field
lower bound. Then port to $p\neq2$: $w_k^h$ solves the \emph{linear}
equation $L_{K_k}w_k^h=\text{RHS-difference}$ with the GJS kernel
$K_k\simeq|v_k(x)-v_k(y)|^{p-2}|x-y|^{-n-sp}$ (exactly the object of
GJS \S7); the datum trick is linear-friendly, and the new constraint
$\gamma<\frac{p-1}{1+sp}$ from (c) should be what makes the kernel-datum
oscillation vanish. The cocycle/second-difference steps (4)--(5) need
rethinking for the vanishing-jet ($\nu=1$) formulation --- but there
$a(h)$ may not even appear, since the CP-style Liouville works with
$v$ itself.}

\begin{remark}[Verification verdict on (c$'$), 2026-06-15]\label{rem:cprime-check}
A careful pass over steps (1)--(5) at $p=2$ finds the argument
\emph{sound in outline}, under the single hypothesis
$\gamma<\frac{1}{1+2s}$ (the tail condition). Specifically:
\begin{itemize}
\item[(1)] $|\nabla H_k|$ splits into an annulus piece
  $\lesssim C_h\rho_k^{\al-2s-2}$ (with $\al-2s-2<0$ since $\al<3<2s+2$
  for $s>\frac12$) and a cap piece $\lesssim\theta_k^{-1}r_k^{2s+1-\al}$
  (with $2s+1-\al>0\iff\gamma<\frac1{1+2s}$); both $\to0$. So
  $H_k=a_k+o(1)$ in $C^1(B_R)$, $a_k$ constant. \checkmark
\item[(2)] \textbf{Not circular.} Decompose
  $(-\Delta)^s\tilde w_k(x_0)=$ (local singular part)$+$(tail). The local
  part is bounded by the uniform interior $[\,\cdot\,]_{C^{2s+\gamma}}$
  bound on $w_k$ (PTU (4.8), $2s+\gamma>2s$); the tail
  $\|\tilde w_k\|_{L^1_s}=\int_{|y|>1}|w_k(y)|\,|y|^{-n-2s}\,dy$ uses
  \emph{only} the growth bound $|w_k(y)|\le C_h|y|^{\al-1}$ on the support
  $B_{\rho_k}$, which converges (uniformly in $k$) iff $\al-1-2s<0$, again
  $\gamma<\frac1{1+2s}$. Neither bound uses $a_k$ or any lower bound on
  $\theta_k$. Then $a_k=(-\Delta)^s\tilde w_k(x_0)-f_k(x_0)-o(1)$ is
  bounded a posteriori. \checkmark \emph{(This was the flagged ``riskiest''
  step; it holds. The truncation manufactures the very $L^1_s$ control PTU
  could not get globally.)}
\item[(3)] Dominated convergence ($|w_k|\le C_h(1+|y|)^{\al-1}\in L^1_s$,
  pointwise $\tilde w_k(y)\to w_0(y)$ since eventually $y\in B_{\rho_k}$)
  gives $\tilde w_k\to w_0$ in $L^1_s$ \emph{globally}; CS stability then
  legitimately applies (its global hypothesis is now met) and yields
  $(-\Delta)^s(v_0(\cdot+h)-v_0)=a(h)$. \checkmark
\item[(4)] Cocycle $w^{h_1+h_2}=w^{h_1}(\cdot+h_2)+w^{h_2}$ plus
  translation invariance of $(-\Delta)^s$ gives $a(h_1{+}h_2)=a(h_1)+a(h_2)$;
  with measurability (pointwise limit of $h\mapsto a_k(h)$, continuous),
  $a(h)=b\cdot h$. \checkmark \emph{Minor: extract one diagonal subsequence
  over a countable dense set of $h$, then use continuity, so $a(\cdot)$ is
  well defined.}
\item[(5)] $\Delta^2_hv_0$ is entire $s$-harmonic with growth
  $(1+|x|)^{\al-2}$, $\al-2<1$, hence \emph{constant} (sublinear-growth
  Liouville); thus $\Delta^3_hv_0\equiv0$, so $v_0$ is a quadratic
  polynomial; the vanishing $2$-jet at the branching point forces
  $v_0\equiv0$, contradicting the normalization. \checkmark
\end{itemize}
\textbf{Conclusion: (c$'$) repairs PTU (4.12) at $p=2$ with no lower bound
on $\theta_k$.} The mechanism is truncate $\to$ oscillation vanishes $\to$
constant bounded a posteriori $\to$ stability now legal $\to$ second
differences annihilate the constant.

\textbf{The $p\neq2$ obstruction, located precisely.} The datum steps
(1)--(3) port (the increment $w^h_k$ solves the linear $L_{K_k}$ equation,
Remark~\ref{rem:linkernel}, and truncation is linear-friendly). But steps
(4)--(5) use \emph{translation invariance} of the operator
($(-\Delta)^sw^{h_1}(\cdot+h_2)=a(h_1)$), and the $v_k$-dependent kernel
$K_k$ is \emph{not} translation invariant. So the cocycle/second-difference
finish does not transplant. The fix is to drop the additive-constant route
altogether and use the vanishing-jet (CP) Liouville on $v_0$ directly,
where $a(h)$ never appears --- consistent with the restatement in
Conjecture~\ref{conj:liouville}. \emph{This is the one item to confirm with
R.T.: that the vanishing-jet formulation removes the need for steps
(4)--(5) in the nonlinear setting.}
\end{remark}

\emph{(d) A warning on pointwise evaluation.} Evaluating
$\fpl{\tilde v_k}$ classically at a point with only $C^{1,\al'}$
regularity needs (after odd-cancellation of the gradient term in the
P.V.) the remainder exponent $p-1+\al'-sp>-... $, i.e.\
$\al'>sp-(p-1)$ when $sp>p-1$ --- which \emph{clashes} with the upper
bound in (a): the window $sp-(p-1)<\al'\le\frac{sp-(p-1)}{p-1}$ is empty
for $p>2$ and borderline-empty at $p=2$. \TODO{This suggests (a) and (d)
do not actually bind simultaneously (PTU work at order $\nu=2$ where the
bookkeeping differs, and weak formulations avoid pointwise evaluation
altogether). Sort out which constraints are real in the weak framework
--- priority for the next session, fresh head needed.}

\subsection*{Synthesis: intersecting the regions (2026-06-15)}
Working at order $\nu=1$ (the only order available for $p\neq2$, since
quadratics are not $(s,p)$-harmonic), the explicit constraints are
\[
  \text{(C1) tail: }\gamma<\tfrac{(p-1)^2}{p-1+sp},\quad
  \text{(C2) jet }\beta<1\text{: }\gamma<(p-1)-\tfrac{sp}{2},\quad
  \text{(C4) datum: }\gamma<\tfrac{p-1}{1+sp},
\]
together with the soft \emph{ceiling} (C3) $\;\al<1+\al_0\;$ and the
\cite{GJS} window $p\in[2,\tfrac{2}{1-s})$ (equivalently $s>1-\tfrac2p$).
Three findings from the bookkeeping (all verified symbolically):
\begin{itemize}
\item[(i)] \textbf{C4 is the binding explicit constraint for $p>2$.} One
  has $C4\le C1$ with equality only at $p=2$
  ($\frac{C4}{C1}=\frac{p-1+sp}{(p-1)(1+sp)}\le1$); and C4 is typically
  below C2 as well (e.g.\ $p=3,s=\tfrac34$: $C4=0.62<C2=0.88<C1=0.94$).
  So the explicit region is $\gamma<\min(C2,C4)$ and is \emph{nonempty}.
\item[(ii)] \textbf{Restricting to $\nu=1$ is not fatal for the new
  ($p>2$) range.} The $\nu=1$ cap C2 reads $\gamma<(p-1)-\tfrac{sp}{2}$;
  at $s=1$ this is $\gamma<\tfrac p2-1$, which is \emph{positive for every
  $p>2$}. Only at $p=2$ does the window pinch to $0$ as $s\nearrow1$ ---
  which is exactly why PTU needed $\nu=2$ there, and exactly the case
  already covered by \cite{PTU}. So the unavailability of $\nu=2$ costs us
  nothing where the result is new.
\item[(iii)] \textbf{The real binding constraint is the ceiling (C3),
  $\al<1+\al_0$.} Since $\al=\frac{sp}{p-1-\gamma}>1$ in the relevant
  range, this needs the \cite{GJS} exponent to satisfy
  $\al_0>\al-1=\frac{sp-(p-1)+\gamma}{p-1-\gamma}$. At $p=3,s=\tfrac34,
  \gamma=\tfrac12$ this demands $\al_0>\tfrac12$ --- a \emph{quantitative}
  lower bound on $\al_0$ that \cite{GJS} (existence of some $\al_0$) does
  not by itself supply.
\end{itemize}
\TODO{Decision point for the meeting. The explicit region
$\gamma<\min(C2,C4)$ is fine; the theorem is conditional on a
\emph{quantitative} $\al_0$ from \cite{GJS} (or a regime where $\al-1$ is
small, i.e.\ $\gamma$ small and $sp$ close to $p-1$). Two honest options:
(A) state the main theorem with hypothesis ``$\al<1+\al_0$'' and carry
$\al_0$ as an abstract constant (clean, but the admissible set is then
only proved nonempty in the perturbative corner $sp\downarrow p-1$);
(B) extract or lower-bound $\al_0(n,s,p)$ from the \cite{GJS} proof
(harder, but yields an explicit region). Ask R.T.\ which framing he
prefers; this is the analogue of PTU's clean $\gamma\in(0,1/3)$.}

%======================================================================
\section{The load-bearing wall: a Liouville theorem}\label{sec:liouville}
%======================================================================
\begin{conjecture}[Liouville with vanishing jet, $\nu=1$;
restated 2026-06-15, cap $\beta<1$ via increments]\label{conj:liouville}
Let $v\in C^{1}_{loc}(\R^n)$ be an entire $(s,p)$-harmonic function with
\[
  v(0)=|Dv(0)|=0
  \qquad\text{and}\qquad
  |v(x)|\le C\,(1+|x|)^{1+\beta}\quad\text{in }\R^n,
\]
for some $0<\beta<1$ (equivalently $\al<2$, i.e.\ Lemma~\ref{lem:beta},
C2), with the tail condition \emph{(C1)} $(\,(1+\beta)(p-1)<sp+(p-1)\,$,
i.e.\ Lemma~\ref{lem:tail}$)$ in force, and assume $v$ is
\emph{nondegenerate} ($Dv\not\equiv0$ on any open set). Then $v\equiv0$.
\end{conjecture}
\begin{remark}[Why the cap is $\beta<1$, not $\beta<\al_0$ --- the route is
forced, 2026-06-15]\label{rem:routeforced}
The earlier ``$\beta<\al_0$'' form copied \cite{CP}'s \emph{local} Lemma~5,
whose proof rescales $v_l=v(l\cdot)/l^{1+\beta}$ and applies the
\emph{sharp} $C^{1,\al_0}$ estimate to $v_l$. \textbf{That route is blocked
nonlocally:} $v_l$ has the same growth $1+\beta$ as $v$, and
$\Tail(v_l)<\infty$ would require $(1+\beta)(p-1)<sp$ --- but the target
$1+\beta=\al=\frac{sp}{p-1-\gamma}$ gives $\al(p-1)-sp=\frac{\gamma ps}{p-1-\gamma}>0$
for \emph{every} $\gamma>0$, so the tail diverges; truncating leaves a
far-field datum with the same divergent exponent, so \cite{GJS} cannot be
applied to $v_l$ at all. The route that \emph{does} survive nonlocally is
PTU's: pass to increments $w^h=v(\cdot+h)-v$, which lower the growth by one
order to $\beta<1$ and (under C1) have finite tail. This is why the correct
cap is $\beta<1$ ($\al<2$), exactly the cap of PTU Theorem~3.1, and why the
sharp exponent $\al_0$ is needed only for compactness (any $\al_0>0$), not
quantitatively. \emph{This supersedes synthesis finding (iii) and the
``need a quantitative $\al_0$'' worry: under the increment route the
binding constraints are C1 and C2, both explicit.} The price is instead the
\textbf{nondegeneracy} hypothesis, which controls the $v$-dependent kernel
$K_v^h$ (Proposition~\ref{prop:linincr}, Remark~\ref{rem:linkernel}).
\end{remark}
\begin{lemma}[Jet-to-growth; \cite{CP} Lemma 6, transplanted]\label{lem:jet}
Let $u\in C^{1+\al'}(B_1)$ with $0<\al'<1$, where
$[u]_{C^{1+\al'}}:=[Du]_{C^{0,\al'}}$. If $u(0)=|Du(0)|=0$ and
\[
  \sup_{0<r<1/2} r^{\al'-\beta}\,[u]_{C^{1+\al'}(B_r)} \le A
  \qquad\text{for some } \beta>\al',
\]
then $|u(x)|\le A|x|^{1+\beta}$ for $x\in B_{1/2}$.
\end{lemma}
\begin{proof}
Fix $x\in B_{1/2}$. Since $u(0)=0$, the mean value theorem gives
$|u(x)|\le|Du(\xi)|\,|x|$ for some $\xi\in B_{|x|}$. Since $Du(0)=0$,
\[
  |Du(\xi)| = |Du(\xi)-Du(0)|
  \le [u]_{C^{1+\al'}(B_{|x|})}\,|\xi|^{\al'}
  \le A\,|x|^{\beta-\al'}\,|x|^{\al'} = A|x|^{\beta},
\]
and the claim follows. (The argument is operator-free: it uses only the
vanishing jet and the seminorm decay.)
\end{proof}
\begin{proof}[Proof strategy (route (b), increments; 2026-06-15)]
The local route (a) of \cite{CP} (rescale $v$, apply the sharp estimate)
is \emph{not available} nonlocally, by Remark~\ref{rem:routeforced}: the
tail of the rescaled $v_l$ diverges for every $\gamma>0$. We therefore
follow PTU's increment route, made linear by
Proposition~\ref{prop:linincr}.
\begin{enumerate}
\item \emph{Increment equation.} For a unit vector $h$ (small), set
  $w^h:=v(\cdot+h)-v$. By Proposition~\ref{prop:linincr} (with $\fpl{v}=0$,
  so the right-hand side is $0$),
  \[
    L_{K_v^h}w^h=0,\qquad
    K_v^h(x,y)=(p-1)\!\int_0^1|(v(x){-}v(y))+t(w^h(x){-}w^h(y))|^{p-2}dt .
  \]
\item \emph{Lowered growth and finite tail.} Since $v\in C^1$ with
  $Dv(0)=0$ and $|v|\le C(1+|x|)^{1+\beta}$, one has
  $|w^h(x)|\le|h|\sup_{[x,x+h]}|Dv|\le C_h(1+|x|)^{\beta}$ with $\beta<1$;
  under (C1) this growth is in the tail space of $L_{K_v^h}$
  (Lemma~\ref{lem:tailfix}).
\item \emph{Linear regularity + decay.} On compact sets where $v$ is
  nondegenerate, $K_v^h$ lies in the measurable-kernel ellipticity class of
  \cite{GJS} \S7, which supplies an interior $C^{1,\al_0}$ estimate with
  tail term for $w^h$. The PTU rescaling/decay argument
  ($w^h_\rho:=\rho^{-\beta}w^h(\rho\cdot)$, interior estimate, then
  $\|Dw^h\|_{L^\infty(B_{1/2\rho})}\le C\rho^{\beta-1}\to0$ as $\rho\to\infty$
  because $\beta<1$) forces $Dw^h(0)=0$; running it at every centre (where
  ellipticity is locally uniform --- this replaces translation invariance)
  gives $Dw^h\equiv0$, so $w^h$ is constant.
\item \emph{Reconstruction.} $w^h\equiv c(h)$ for all small $h$; the cocycle
  $w^{h_1+h_2}=w^{h_1}(\cdot+h_2)+w^{h_2}$ forces $c$ additive and (by
  continuity) linear, $c(h)=b\cdot h$; but $w^h(0)=v(h)$ has a double zero
  at $h=0$, so $b=0$ and $v$ is constant on the lattice of increments,
  hence (with the vanishing jet) $v\equiv0$.
\end{enumerate}
\TODO{Two genuine gaps, both to raise with R.T.: (i) \textbf{nondegeneracy}
--- that the blow-up limit $v_0$ has $Dv_0\not\equiv0$ on open sets, so
$K_{v_0}^h$ is uniformly elliptic on compacta (near the branching point $v$
is flat, but that is a single point); (ii) the precise \textbf{linear
nonlocal Liouville} ``$L_K w=0$, $K$ measurable elliptic, sublinear growth
$\Rightarrow$ constant'' --- locate the exact reference (Kassmann; or
\cite{GJS} \S7 run as a Liouville) with the tail hypothesis matching (C1).
Step (3)'s ``run at every centre'' needs the ellipticity constants locally
uniform, which is where (i) enters.}
\end{proof}
\begin{lemma}[Resolution of the nonlocal tail gap via increments;
2026-06-15]\label{lem:tailfix}
Let $v$ satisfy the hypotheses of Conjecture~\ref{conj:liouville}, so
$|v(x)|\le C(1+|x|)^{1+\beta}$ with $\beta=\al-1$. The rescaled
$v_l:=v(l\cdot)/l^{1+\beta}$ has \emph{divergent} tail
$\Tail(v_l)=\infty$, because $(1+\beta)(p-1)=\al(p-1)>sp$; thus
\cite{GJS} cannot be applied to $v_l$ directly. Pass instead to first
increments $w^h:=v(\cdot+h)-v(\cdot)$. Since $v\in C^1$ with $Dv(0)=0$,
\[
  |w^h(x)|\le |h|\sup_{[x,x+h]}|Dv|\;\lesssim\;|h|\,(1+|x|)^{\beta},
\]
one order lower, and the nonlocal tail of $w^h$ is finite iff
\[
  \beta(p-1)=(\al-1)(p-1)<sp,
  \qquad\text{i.e.\ iff the tail condition \emph{(C1)}}\
  \gamma<\tfrac{(p-1)^2}{p-1+sp}.
\]
Hence the incremental quotient — not $v_l$ — is the object on which the
scaling/decay argument can run.
\end{lemma}
\begin{proposition}[Increments solve a linear nonlocal equation;
2026-06-15]\label{prop:linincr}
Let $v\in C^1_{loc}\cap L^\infty_{loc}$ and $h\in\R^n$, and set
$w^h:=v(\cdot+h)-v(\cdot)$. Then, wherever the principal values exist,
\[
  \fpl{\big[v(\cdot+h)\big]}(x)-\fpl{v}(x)
  = L_{K_v^h}w^h(x)
  := c_{n,s,p}\,\mathrm{P.V.}\!\int_{\R^n}
     \big(w^h(x)-w^h(y)\big)\,K_v^h(x,y)\,\frac{dy}{|x-y|^{n+sp}},
\]
with the symmetric, \emph{strictly positive} kernel
\[
  K_v^h(x,y):=(p-1)\int_0^1
   \big|\,(v(x)-v(y))+t\,(w^h(x)-w^h(y))\,\big|^{\,p-2}\,dt\;\ge0 .
\]
Consequently, if $v=v_k$ is a blow-up solution as in
Proposition~\ref{prop:homog}, then
\[
  L_{K_{v_k}^h}w_k^h
  = \theta_k^{\,\gamma-(p-1)}\,
    \Big[\big(u\mapsto u_+^\gamma-u_-^\gamma\big)(v_k)(\cdot+h)
         -\big(u\mapsto u_+^\gamma-u_-^\gamma\big)(v_k)\Big]
  =: g_k^h,
\]
and $g_k^h\to0$ locally uniformly. So $w_k^h$ solves a \emph{linear}
nonlocal equation with right-hand side tending to $0$ --- the structure
that makes the datum/truncation argument of \textnormal{(c$'$)} available
for $p\neq2$.
\end{proposition}
\begin{proof}
Translating $y\mapsto y+h$ in the integral for
$\fpl{[v(\cdot+h)]}(x)$ writes it as
$c\,\mathrm{P.V.}\!\int \Jp{\big(V(x)-V(y)\big)}|x-y|^{-n-sp}dy$ with
$V:=v(\cdot+h)$. Subtract $\fpl{v}(x)$ and use, for the scalar map
$J_p(t)=\Jp{t}$ with $J_p'(t)=(p-1)|t|^{p-2}$ (continuous for $p\ge2$),
the exact identity
\[
  J_p(a+\delta)-J_p(a)
  =\Big((p-1)\!\int_0^1|a+t\delta|^{p-2}\,dt\Big)\,\delta,
\]
applied pointwise with $a=v(x)-v(y)$ and
$\delta=w^h(x)-w^h(y)$ (so $a+\delta=V(x)-V(y)$). The bracket is
$K_v^h(x,y)$; symmetry in $(x,y)$ and nonnegativity are immediate. The
displayed right-hand side for $w_k^h$ is the increment of
$\fpl{v_k}=\theta_k^{\gamma-(p-1)}[(v_k)_+^\gamma-(v_k)_-^\gamma]$ from
Proposition~\ref{prop:homog}, and the prefactor $\to0$. \emph{(Identity
verified numerically for $p\in\{2,2.5,3,4\}$, $K_v^h>0$.)}
\end{proof}
\begin{remark}[What changes for $p\neq2$]\label{rem:linkernel}
At $p=2$, $w^h$ is again $s$-harmonic (linearity) and
PTU Theorem~3.1 finishes. For $p\neq2$, $w^h$ solves not the homogeneous
equation but a \emph{linear} nonlocal equation $L_{K_v}w^h=0$ with the
$v$-dependent kernel
$K_v(x,y)\simeq|v(x)-v(y)|^{p-2}|x-y|^{-n-sp}$ (the linearization object of
\cite{GJS} \S7). Two consequences. (i) The $C^{\al_0}$/$C^{1,\al_0}$
estimate one applies is the one for this \emph{measurable-kernel} linear
class (\cite{GJS} \S7), not Theorem~1.1 for $\fpl{w}=0$; its tail term is
finite exactly under (C1). (ii) \textbf{Caveat (degeneracy):} $K_v$
degenerates where $v$ is flat ($|v(x)-v(y)|$ small). At the branching
point $v$ has a vanishing $1$-jet, so the kernel is weakest precisely
there; one must check the blow-up limit $v_0$ is not flat on an open set
(its gradient does not vanish identically away from $0$) so that $K_{v_0}$
lies in a uniform ellipticity class on the relevant region. This is the
genuine technical content of the $p\neq2$ Liouville and should be flagged
to R.T.
\end{remark}
\begin{remark}
$\nu=2$ is out of reach: quadratic polynomials are not $(s,p)$-harmonic
for $p\neq2$, so even the statement fails at order two; consistent with
the absence of $C^2$ theory for the (local or nonlocal) $p$-Laplacian.
State this as a remark in the paper, mirroring PTU Remark~4.1 in reverse.
\end{remark}
%======================================================================
\section{Target theorems}\label{sec:targets}
%======================================================================
\begin{theorem}[Main growth estimate --- target]\label{thm:main}
Let $(s,p,\gamma)$ lie in the admissible region of
\S\ref{sec:constraints} and let $u\in L^\infty(\R^n)$ solve
\eqref{eq:main}. If $x_0\in B_{1/2}$ satisfies $u(x_0)=|Du(x_0)|=0$, then
\[
  |u(x)| \le C\,\|u\|_{L^\infty(\R^n)}\;|x-x_0|^{\frac{sp}{p-1-\gamma}},
  \qquad x\in B_{1/2}(x_0),
\]
with $C=C(n,s,p,\gamma)>0$.
\end{theorem}
\begin{proof}[Proof modulo Conjecture~\ref{conj:liouville}]
We run the device of \cite{CP} \S6 (the quasilinear form of PTU
(4.3)--(4.10)) with our ingredients. \emph{Choice of working exponent.}
Pick any $\al'$ with
\[
  0<\al'<\min\{\al_0,\ \al-1\}, \qquad \al=\tfrac{sp}{p-1-\gamma},
\]
a nonempty range because $\al>1$ on the admissible set and
$\al_0>0$ (Lemma~\ref{lem:ceiling}). \emph{This $\al'$ is generic and
sub-sharp; the sharp $\al_0$ is used only for this compactness step, not
quantitatively.} We may take $x_0=0$, $\|u\|_{L^\infty(\R^n)}=1$, and must show
$|u(x)|\le C|x|^{\al}$.

\emph{Contradiction and the $\theta_k$ device.} If not, there are solutions
$u_k$ of \eqref{eq:main} with $0$ a branching point,
$[u_k]_{C^{1+\al'}(B_{1/2})}\le 2C$, and points $x_k$ with
$|u_k(x_k)|>k|x_k|^{\al}$. Put
\[
  \theta_k(r'):=\sup_{r'<r<1/2} r^{\,1+\al'-\al}\,[u_k]_{C^{1+\al'}(B_r)},
  \qquad 1+\al'-\al<0 .
\]
By Lemma~\ref{lem:jet} in contrapositive form (its hypothesis needs
$\al-1>\al'$, which holds by the choice of $\al'$): were $\theta_k(0^+)\le
k$, the vanishing jet would force $|u_k(x)|\le k|x|^{\al}$, contradicting
$x_k$. Hence $\lim_{r'\to0}\theta_k(r')>k$, and we may pick $r_k>1/k$ with
\[
  r_k^{\,1+\al'-\al}[u_k]_{C^{1+\al'}(B_{r_k})}
  \ge\tfrac1k\theta_k(1/k)\ge\tfrac12\theta_k(r_k),
\]
so $r_k\to0$. Set $v_k(x):=u_k(r_kx)/\big(\theta_k(r_k)\,r_k^{\al}\big)$.
\emph{Seminorm bound.} As in \cite{CP} (33), for $1\le R\le\frac1{2r_k}$,
\[
  [v_k]_{C^{1+\al'}(B_R)}
  =\frac{(r_kR)^{1+\al'-\al}[u_k]_{C^{1+\al'}(B_{r_kR})}}{\theta_k(r_k)}
   \,R^{\al-1-\al'}
  \le R^{\al-1-\al'},
\]
using $(r_kR)^{1+\al'-\al}[u_k]_{C^{1+\al'}(B_{r_kR})}\le\theta_k(r_kR)\le
\theta_k(r_k)$. \emph{Local growth via the gradient anchor.} With
$\eta\equiv1$ on $B_{1/2}$, $\eta=0$ off $B_1$,
$\int_{B_1}\eta\,D_ev_k=\int_{B_1}D_e\eta\,v_k\le C_n$ gives a point $z\in
B_1$ with $|Dv_k(z)|\le C_n$; combined with the seminorm bound,
$|Dv_k(x)|\le CR^{\al-1}$ and, since $v_k(0)=0$, $|v_k(x)|\le C|x|^{\al}$
for $1\le|x|\le\frac1{2r_k}$. \emph{Compactness.} Thus $v_k$ is locally
bounded with uniform $C^{1+\al'}$ bounds, so along a subsequence $v_k\to
v_0$ in $C^{1+\al'}_{loc}(\R^n)$, with
\[
  [v_0]_{C^{1+\al'}(B_R)}\le R^{\al-1-\al'}\ (R\ge1),\qquad
  |v_0(x)|\le C(1+|x|)^{\al},\qquad
  [v_0]_{C^{1+\al'}(B_1)}\ge\tfrac12 .
\]
Since $0$ is a branching point of each $u_k$, $v_0(0)=|Dv_0(0)|=0$. By
Proposition~\ref{prop:homog}, $\fpl{v_0}=0$ in $\R^n$. Its growth exponent
$\al$ satisfies $\al<2$ (Lemma~\ref{lem:beta}, C2), i.e.\ $\beta=\al-1<1$,
so --- provided $v_0$ is nondegenerate --- the hypothesis of
Conjecture~\ref{conj:liouville} holds and $v_0\equiv0$, contradicting
$[v_0]_{C^{1+\al'}(B_1)}\ge\frac12$.
\end{proof}
\NOTE{\textbf{On the exponent bookkeeping (settled 2026-06-15; this note
records the final state after two corrections).} (a) The device runs at a
\emph{generic, sub-sharp} $\al'<\min\{\al_0,\al-1\}$, \emph{not} at level
$1+\al_0$; so the earlier ``$1+\al_0<\al<2$ contradicts the conjecture''
panic was a misidentification --- there is no contradiction, and $\al_0$
enters only for compactness (any $\al_0>0$). (b) But the \emph{opposite}
over-correction (``so the binding constraint is the sharp $\al<1+\al_0$,
CP-style'') is \emph{also} wrong nonlocally: CP's one-shot rescaling is
\textbf{blocked by the tail} (Remark~\ref{rem:routeforced}, verified ---
$\al(p-1)-sp=\frac{\gamma ps}{p-1-\gamma}>0$ always). The Liouville must go
through \emph{increments}, whose cap is $\beta<1$, i.e.\ $\boxed{\al<2}$
(Lemma~\ref{lem:beta}, C2). \textbf{Net: the binding constraints are C1
(tail) and C2 ($\al<2$), both explicit; no quantitative lower bound on
$\al_0$ is needed.} The price moves elsewhere: the increment route needs
the \textbf{nondegeneracy} of $v_0$ (Conjecture~\ref{conj:liouville},
Remark~\ref{rem:routeforced}) and a linear nonlocal Liouville. \emph{This
supersedes synthesis finding (iii).} $\,$Update \S\ref{sec:constraints}
accordingly when the dust settles; for now the corrected story lives here
and in Remark~\ref{rem:routeforced}.}
\begin{corollary}[Sharp growth of the gradient --- target]
Under the conditions of Theorem~\ref{thm:main},
$\sup_{B_r(x_0)}|Du|\le C\,r^{\frac{sp}{p-1-\gamma}-1}$ for $r<1/2$.
\end{corollary}
\begin{theorem}[Two-phase local $p$-dead-core --- target, new even locally]\label{thm:local}
Let $u$ solve $\Delta_p u = u_+^{\gamma}-u_-^{\gamma}$ in $B_1$ with
$(p,\gamma)$ admissible. At points with $u(x_0)=|Du(x_0)|=0$,
\[
  |u(x)|\le C\,|x-x_0|^{\frac{p}{p-1-\gamma}} .
\]
\end{theorem}
\begin{proof}[Proof skeleton]
\TODO{$s\nearrow1$ in Theorem~\ref{thm:main}, after checking
(i) $\fpl{u}\to-\Delta_p u$ in the appropriate sense (locate the
Bourgain--Brezis--Mironescu-type statement for $\fpl{}$ and fix
normalizing constants $c_{n,s,p}$), and
(ii) uniformity in $s$ of every constant in the chain --- in particular
of the $C^{1,\al_0}$ estimate as $s\nearrow1$ (the window
$p<\frac{2}{1-s}$ opens to all $p\ge2$). The local two-phase result does
not appear in da Silva--Rossi--Salort / da Silva--Salort / Teixeira, which
are one-phase: highlight this.}
\end{proof}
\begin{theorem}[Liouville with growth --- target]
Entire bounded solutions of \eqref{eq:main} in $\R^n$ with
$u(x)=o\big(|x-x_0|^{\frac{sp}{p-1-\gamma}}\big)$ as $|x|\to\infty$, for
some $x_0\in\partial\{u>0\}\cap\{|Du|=0\}$, vanish identically.
\end{theorem}
\begin{proof}[Proof skeleton]
\TODO{Transplant PTU Theorem~5.3 verbatim: rescale, apply
Theorem~\ref{thm:main}, squeeze. Should be a one-page corollary once the
main theorem holds.}
\end{proof}
\begin{question}
Nondegeneracy: $\sup_{B_r(x_0)}u_+\ge c\,r^{\frac{sp}{p-1-\gamma}}$ at
free boundary points? Requires a barrier: compute (or bound) $\fpl{}$ of
$x\mapsto |x|^{\al}$. Power-function computations for $\fpl{}$ exist in
the literature (1-d computations of Brasco et al.; $d^s$-barriers of
Iannizzotto--Mosconi--Squassina). Separate note.
\end{question}
%======================================================================
\section{Reading list with shopping lists}
%======================================================================
\begin{itemize}
  \item \textbf{PTU} (this paper, arXiv:2504.09557): \S4 line-by-line;
        extract the $\theta_k$ device and the exponent bookkeeping.
  \item \textbf{Brasco--Lindgren--Schikorra} (arXiv:1711.09835):
        shopping list --- interior $C^{0,\gamma_0}$ estimate with explicit
        tail dependence; the form of $\Tail$; stability under limits.
  \item \textbf{Giovagnoli--Jesus--Silvestre} \cite{GJS}
        (arXiv:2509.26565): \emph{the} $C^{1,\al_0}$ paper for
        Lemma~\ref{lem:ceiling} --- Theorem 1.1 gives interior $C^{1,\al}$
        for $\fpl{u}=0$, $p\in[2,\frac{2}{1-s})$, with tail term.
        Shopping list: constants, uniformity in $s$;
        \textbf{check extension to $\fpl{u}=f\in L^\infty$}
        (their statement is homogeneous).
  \item \textbf{Correa--dos Prazeres} \cite{CP} (arXiv:2504.11370):
        two-phase quenching-type (Alt--Phillips) problem for the local
        $p$-Laplacian; \S6 runs the Liouville + known-regularity strategy
        in the local quasilinear setting --- the closest blueprint for
        \S\ref{sec:liouville}, route (a). [Supervisor's pointer.]
  \item \textbf{Korvenp\"a\"a--Kuusi--Lindgren}: weak $\equiv$ viscosity;
        cite, do not read deeply on the first pass.
  \item \textbf{da Silva--Salort} (arXiv:2501.13063): the local benchmark
        $C^{\frac{p}{p-1-q}}$; compare hypotheses and the one-phase vs
        two-phase gap.
  \item \textbf{Teymurazyan}, \emph{The fractional Laplacian: a primer}:
        background reference; the author is at KAUST --- bring \S3--\S5 of
        these notes to his office once Lemmas~\ref{lem:tail}--\ref{lem:ceiling}
        are verified.
\end{itemize}
\begin{thebibliography}{9}
\bibitem{PTU} D. dos Prazeres, R. Teymurazyan, J.M. Urbano,
\emph{Improved regularity for a nonlocal dead-core problem},
arXiv:2504.09557, to appear in Indiana Univ. Math. J.
\bibitem{CS-approx} L. Caffarelli, L. Silvestre,
\emph{Regularity results for nonlocal equations by approximation},
Arch. Ration. Mech. Anal. 200 (2011), 59--88.
\bibitem{GJS} D. Giovagnoli, D. Jesus, L. Silvestre,
\emph{$C^{1+\alpha}$ regularity for fractional $p$-harmonic functions},
arXiv:2509.26565.
\bibitem{CP} J.C. Correa, D. dos Prazeres,
\emph{A two-phase quenching-type problem for the $p$-Laplacian},
arXiv:2504.11370.
\end{thebibliography}
\end{document}
